\chapter{INTRODUÇÃO}
\label{capitulo-introducao-proposta}

Esta proposta de extensão propõe o desenvolvimento e a avaliação de um protótipo de bracelete inteligente baseado na plataforma microcontrolada ESP32, concebido como tecnologia assistiva de baixo custo voltada a auxiliar pessoas com baixa visão na percepção espacial e orientação diante de obstáculos físicos. A locomoção independente de pessoas com deficiência visual enfrenta frequentes barreiras em ambientes urbanos e institucionais. Enquanto a tradicional bengala longa é eficiente no mapeamento tátil do solo e identificação de degraus, obstáculos aéreos e suspensos --- como ramagens, placas de sinalização, orelhões, caçambas e móveis --- situam-se acima da linha de varredura, expondo os usuários a riscos contínuos de colisão.

O tema se faz oportuno para a Ciência da Computação ao articular de forma direta os conhecimentos de sistemas embarcados, programação de baixo nível em C/C++, processamento de sinais de sensores de proximidade e interação humano-dispositivo com uma necessidade de caráter social. O projeto engloba valores de acessibilidade, cidadania e autonomia, alinhando-se aos princípios da Lei Brasileira de Inclusão da Pessoa com Deficiência (Lei nº 13.146/2015). Para a formação dos discentes, a proposta proporciona a vivência prática de conceber, montar e testar uma solução tecnológica completa de engenharia de software e hardware, respondendo a demandas concretas da comunidade.

\section{Identificação da Ação Extensionista}
\label{sec:identificacao-acao-proposta}

Em conformidade com as diretrizes do modelo de proposta de extensão do IFPR, detalham-se a seguir as características essenciais da ação:

\begin{itemize}
    \item \textbf{Público-alvo:} Pessoas com baixa visão e deficiência visual que necessitam de auxílio para a detecção antecipada de obstáculos suspensos, bem como a comunidade técnica e escolar interessada em tecnologias assistivas abertas e acessíveis;
    \item \textbf{Quando ocorrerá:} A execução do projeto está prevista para o primeiro semestre letivo de 2026, com cronograma concentrado ao longo de quatro semanas entre os meses de março e abril de 2026, com atividades laboratoriais semanais;
    \item \textbf{Onde:} As atividades de projeto, montagem, programação e testes controlados serão realizadas no Instituto Federal do Paraná (IFPR) -- Campus Pinhais, especificamente nas dependências dos Laboratórios de Informática e Sistemas Embarcados;
    \item \textbf{Quem são os extensionistas:} A execução da proposta será conduzida pelos estudantes Matheus Vitor Lourenço Schionato e Vairtles Nehisen Mounkassa Liel, discentes do curso de Bacharelado em Ciência da Computação, sob a orientação e o acompanhamento do professor responsável pela disciplina de Práticas de Extensão III.
\end{itemize}
